% ===== PRÉAMBULE CANONIQUE — à coller VERBATIM en tête de CHAQUE .tex =====
\documentclass[11pt,a4paper]{article}
\usepackage[utf8]{inputenc}\usepackage[T1]{fontenc}\usepackage{lmodern}
\usepackage[french]{babel}
\usepackage{amsmath,amssymb,mathtools}\usepackage{icomma}
\usepackage{siunitx}\sisetup{locale=FR}  % \qty{}{}, \si{}, \ang{}
\usepackage{xcolor}\usepackage{colortbl}
\usepackage{tikz}\usepackage{pgfplots}\pgfplotsset{compat=1.18}
\usepackage[siunitx,europeanresistors]{circuitikz} % circuits (élec.)
\usepackage[most]{tcolorbox}\usepackage{enumitem}\usepackage{array,booktabs}
\usepackage[a4paper,margin=2.2cm]{geometry}
\usetikzlibrary{arrows.meta,calc,angles,quotes,positioning,patterns,%
  backgrounds,shapes.geometric,decorations.pathmorphing,decorations.markings,intersections}
\definecolor{primary}{RGB}{222,49,99}\definecolor{primarydark}{RGB}{150,22,55}
\definecolor{defcol}{RGB}{0,114,178}\definecolor{methcol}{RGB}{52,92,148}
\definecolor{excol}{RGB}{0,135,90}\definecolor{attncol}{RGB}{200,40,55}
\tcbset{boxbase/.style={enhanced,breakable,boxrule=0.9pt,arc=2.5pt,
  left=3mm,right=3mm,top=2.3mm,bottom=2.3mm,fonttitle=\bfseries,coltitle=white}}
% --- boîtes TUTORIEL (noms EXACTS, ne pas renommer) ---
\newtcolorbox{cle}[1][]{boxbase,colback=defcol!6,colframe=defcol,
  title={\textsc{À retenir}\ifx\relax#1\relax\relax\else\textmd{ : \itshape #1}\fi}}
\newtcolorbox{methode}[1][]{boxbase,colback=methcol!7,colframe=methcol,sharp corners,
  title={\textsc{Méthode}\ifx\relax#1\relax\relax\else\textmd{ --- \itshape #1}\fi}}
\newtcolorbox{exemplebox}[1][]{boxbase,colback=excol!5,colframe=excol,
  title={\textsc{Exemple}\ifx\relax#1\relax\relax\else\textmd{ : \itshape #1}\fi}}
\newtcolorbox{piege}[1][]{boxbase,colback=attncol!6,colframe=attncol,
  title={\textsc{Piège}\ifx\relax#1\relax\relax\else\textmd{ : \itshape #1}\fi}}
\newtcolorbox{remarque}[1][]{boxbase,colback=black!4,colframe=black!45,
  title={\textsc{Remarque}\ifx\relax#1\relax\relax\else\textmd{ : \itshape #1}\fi}}
% --- boîtes EXERCICE / CORRIGÉ (noms EXACTS, auto-comptées) ---
\newtcolorbox[auto counter]{exo}[1][]{boxbase,colback=primary!4,colframe=primary,
  title={\textsc{Exercice}~\thetcbcounter\ifx\relax#1\relax\relax\else\textmd{ --- \itshape #1}\fi}}
\newtcolorbox[auto counter]{corrige}[1][]{boxbase,colback=excol!4,colframe=excol,
  title={\textsc{Corrigé}~\thetcbcounter\ifx\relax#1\relax\relax\else\textmd{ --- \itshape #1}\fi}}
\newcommand{\vv}[1]{\vec{#1}}\newcommand{\ds}{\displaystyle}
\newcommand{\R}{\mathbb{R}}\newcommand{\N}{\mathbb{N}}\newcommand{\Z}{\mathbb{Z}}
% =========================================================================
\begin{document}
\usetikzlibrary{babel}
\begin{exo}[trois lampes : vrai ou faux ?]
Trois lampes \emph{identiques} (assimilées à des résistances égales) sont montées ainsi : $L_1$ est
\textbf{en série} avec le groupe formé de $L_2$ et $L_3$ \textbf{en parallèle}. Le générateur débite un
courant total $I$.
\begin{center}
\begin{circuitikz}[scale=0.95,transform shape,>=Stealth]
  \coordinate (A) at (2.4,0); \coordinate (B) at (5.2,0);
  \draw (0,-1.7) to[battery1,l_=$G$] (0,0) to[lamp,l=$L_1$] (A);
  \draw (A) -- (2.4,0.95) to[lamp,l=$L_2$] (5.2,0.95) -- (B);
  \draw (A) -- (2.4,-0.95) to[lamp,l=$L_3$] (5.2,-0.95) -- (B);
  \draw (B) -- (6.4,0) -- (6.4,-1.7) -- (0,-1.7);
  \draw[->,primary,line width=1pt] (0.55,0.28) -- (1.05,0.28) node[above,font=\scriptsize,primary]{$I$};
  \draw[->,primary] (2.95,0.95) -- (3.45,0.95) node[above,font=\scriptsize,primary]{$I_2$};
  \draw[->,primary] (2.95,-0.95) -- (3.45,-0.95) node[below,font=\scriptsize,primary]{$I_3$};
  \fill (A) circle(1.8pt); \fill (B) circle(1.8pt);
  \node[above,font=\footnotesize] at (2.4,0.08){$A$};
\end{circuitikz}
\end{center}
Réponds par \textbf{vrai} ou \textbf{faux}, en justifiant :
\begin{enumerate}[label=\alph*)]
  \item l'intensité dans $L_2$ est égale à celle dans $L_3$ ;
  \item l'intensité qui traverse $L_1$ est le \emph{double} de celle qui traverse $L_2$ ;
  \item un ampèremètre placé dans la branche de $L_1$ et un ampèremètre placé juste à la sortie du
        générateur indiquent des valeurs différentes ;
  \item si $I=\qty{0.6}{\ampere}$, alors $L_2$ et $L_3$ sont traversées chacune par $\qty{0.3}{\ampere}$.
\end{enumerate}
\end{exo}

\begin{exo}[deux nœuds en cascade]
Dans la portion de circuit ci-dessous, un courant $I=\qty{8}{\ampere}$ arrive au nœud $N_1$. De $N_1$
partent $I_a$ (vers $N_2$) et $I_b=\qty{3}{\ampere}$. Au nœud $N_2$ arrivent $I_a$ et un courant
$I_c=\qty{2}{\ampere}$ ; il en part $I_d$.
\begin{center}
\begin{circuitikz}[scale=1,transform shape,>=Stealth]
  \coordinate (N1) at (0,0); \coordinate (N2) at (3.2,0);
  \draw[->,primary,line width=1pt] (-1.7,0) -- (-0.15,0) node[midway,above,font=\footnotesize,primary]{$I$};
  \draw[thick] (N1) -- (N2); \draw[->,primary] (1.1,0) -- (1.7,0) node[above,font=\scriptsize,primary]{$I_a$};
  \draw[thick] (N1) -- (0,-1.4); \draw[->,primary] (0,-0.55) -- (0,-1.05) node[right,font=\scriptsize,primary]{$I_b$};
  \draw[thick] (3.2,1.4) -- (N2); \draw[->,primary] (3.2,1.05) -- (3.2,0.55) node[right,font=\scriptsize,primary]{$I_c$};
  \draw[thick] (N2) -- (4.9,0); \draw[->,primary] (3.9,0) -- (4.5,0) node[above,font=\scriptsize,primary]{$I_d$};
  \fill (N1) circle(2pt); \node[above,font=\footnotesize] at (0,0.12){$N_1$};
  \fill (N2) circle(2pt); \node[below,font=\footnotesize] at (3.2,-0.12){$N_2$};
\end{circuitikz}
\end{center}
Calcule $I_a$ puis $I_d$.
\end{exo}

\begin{exo}[comptage d'électrons — style concours]
Un fil conducteur de section $\qty{4}{\milli\metre\squared}$ est parcouru par un courant
$I=\qty{3.2}{\ampere}$. On donne $e=\qty{1.6e-19}{\coulomb}$ et $N_A=6{,}02\times10^{23}$.
Après avoir calculé le nombre $n$ d'électrons traversant une section \emph{chaque seconde}, indique
si chacune de ces affirmations est vraie ou fausse :
\begin{enumerate}[label=\alph*)]
  \item $n$ est supérieur à une mole d'électrons ;
  \item $n$ est inférieur à $10^{13}$ électrons ;
  \item $n$ est supérieur à $1{,}8\times10^{13}$ électrons.
\end{enumerate}
\end{exo}

\begin{exo}[l'éclair : un courant intense et bref]
Un éclair transporte une charge $Q=\qty{15}{\coulomb}$ en une durée $t=\qty{2}{\milli\second}$.
On donne $e=\qty{1.6e-19}{\coulomb}$.
\begin{enumerate}[label=\alph*)]
  \item Calcule l'intensité moyenne du courant de l'éclair.
  \item Combien d'électrons cette charge représente-t-elle ?
  \item Compare cette intensité à celle d'un circuit domestique ($\approx\qty{10}{\ampere}$).
\end{enumerate}
\end{exo}

\end{document}
